%% ============================================================
%% CHAPTER 1 — Introduction            (budget: ~4,000 words)
%% Role: state the question, the claim, and the map of the thesis.
%% ============================================================
\chapter{Introduction}
\label{ch:intro}

\section{Motivation}
\label{sec:intro-motivation}
% Para 1: LLM agents moved from chatbots to systems that act — the ReAct
%   loop as the unit of agency. One paragraph, cite \citep{yao2023react}.
% Para 2: the runtime shrink history: OpenClaw (400k+ LoC TS) -> nanobot
%   (4k Python) -> PicoClaw/ZeroClaw (<10MB Go/Rust). Each step disproved a
%   dependency assumed essential. All stop at Linux.
% Para 3: why the OS boundary matters: cost (a $4 device vs a $50 SBC),
%   power (always-on economics), deployment surface (billions of MCUs),
%   and intellectually — what IS the minimal sufficient agent runtime?
% Para 4: the enabling observation: intelligence in the cloud, architecture
%   on the device. If reasoning consumes no local resources, the question
%   sharpens to the architectural minimum.

\section{Problem Statement and Research Questions}
\label{sec:intro-rqs}
% One paragraph framing, then the RQs as an enumerated list. Each RQ maps
% to exactly one results chapter — say so explicitly here.
% RQ1 (Ch3+5): What does a complete agent runtime require on bare metal,
%   and at what resource cost? Can it reach task parity with Linux runtimes?
% RQ2 (Ch6): How should the per-step context of a ReAct loop be managed on
%   a device that must re-upload it on every call?
% RQ3 (Ch7): How can an OS-less agent extend its own capabilities at
%   runtime, and what replaces the OS services that skills silently assume?
% RQ4 (Ch8): Can multi-step cloud plans be executed safely on-device to
%   reduce per-step cloud dependence, given irreversible physical actions?

\section{Contributions}
\label{sec:intro-contributions}
% Numbered list, one sentence each, each tied to an RQ and a chapter:
% C1: MimiClaw — the first fully documented bare-metal agent runtime
%     (system + open-source artifact).
% C2: A five-baseline evaluation methodology and measurement study
%     (parity, latency waterfall, energy, fault recovery).
% C3: Identification of within-turn quadratic re-upload + the context
%     compaction ladder (L0–L3) with measured trade-offs.
% C4: A design space and implementation of skill systems without an OS
%     (prompt / declarative / interpreted script), incl. LLM authorability.
% C5: Guarded speculative execution: commit-after-verify semantics for
%     irreversible actions, the silent-divergence failure mode, and its
%     measured call/energy savings.

\section{Publications and Software}
\label{sec:intro-publications}
% Placeholder: list arXiv reports / workshop papers as they appear, and the
% open-source repositories (mimiclaw mainline, research fork), stating which
% thesis chapters draw on which artifact.

\section{Thesis Outline}
\label{sec:intro-outline}
% One short paragraph per chapter (2–3 sentences), so an examiner can
% navigate. End with a reading-dependency note: Ch3–4 are prerequisites for
% Ch5–8; Ch6–8 are mutually independent.
